

% ACM Large 1-Column Format
\settopmatter{printacmref=false}
\renewcommand\footnotetextcopyrightpermission[1]{}
\pagestyle{plain}

%------------------------------------------------------------------
\section{Variant A - Classification Only Baseline}
%------------------------------------------------------------------

\subsection{The Core Idea}

Imagine you are learning to judge whether two statements agree with
each other, contradict each other, or are simply unrelated. Your
teacher hands you thousands of examples and asks you to answer after
reading each one:

\begin{quote}
\textit{Statement 1: ``A dog is running in the park.''} \\
\textit{Statement 2: ``An animal is outside.''} \\
\textit{Answer: They agree.}
\end{quote}

There is no explanation given. No reasoning. Just the two statements
and the correct label. You are expected to figure out the pattern on
your own.

This is exactly what Variant A does. It is the most direct possible
approach to Natural Language Inference: read the premise, read the
hypothesis, predict the label. Nothing more. The model receives no
guidance about \textit{why} a label is correct. It must discover the
relationship between the two sentences entirely from the training
signal --- the feedback it receives when its prediction is right or
wrong.

This simplicity is precisely what makes Variant A valuable. It is not
trying to be clever. It is trying to establish an honest, reproducible
answer to a foundational question: \textit{how well can a strong
pretrained encoder solve NLI when given nothing but the raw input?}
Every other variant in this project --- B, C, and D --- adds some form
of reasoning or explanation to the training process. Variant A tells us
what we get without any of that. It is the ground truth against which
all improvements must be measured.

%------------------------------------------------------------------
\subsection{Why This Is Harder Than It Looks}
%------------------------------------------------------------------

At first glance, NLI might seem like a simple classification problem.
You have two sentences and three possible labels. But the task is
deceptively difficult. Consider the following pair:

\begin{quote}
\textit{Premise: ``The lawyer argued that the contract was invalid.''} \\
\textit{Hypothesis: ``Someone made a legal claim.''} \\
\textit{Label: Entailment}
\end{quote}

To get this right, the model must understand that a lawyer is a legal
professional, that arguing in a legal context constitutes making a
claim, and that the act of arguing about a contract falls within the
domain of legal activity. None of this is explicit in the text. It
requires world knowledge, semantic inference, and category reasoning.

A model that merely memorizes surface patterns --- that ``lawyer'' and
``legal'' tend to appear together, or that certain sentence structures
imply entailment --- will perform well on examples that resemble its
training data but fail badly the moment it encounters phrasing it has
not seen before. This is the fundamental weakness of
classification-only approaches, and it is what Variant A's
cross-domain evaluation is designed to reveal.

%------------------------------------------------------------------
\subsection{Model Architecture}
%------------------------------------------------------------------

Variant A uses \texttt{microsoft/deberta-v3-base} as its backbone, a
184-million parameter transformer encoder developed by Microsoft.
DeBERTa (Decoding-enhanced BERT with Disentangled Attention) improves
upon BERT and RoBERTa through two key innovations.

First, it uses disentangled attention, where each token is represented
by two separate vectors --- one for its content and one for its
position --- and attention weights are computed using all four
combinations of content and position interactions. This allows the
model to capture relative positional relationships more precisely than
standard attention mechanisms.

Second, it uses an enhanced mask decoder during pre-training that
reintroduces absolute position information at the output layer, helping
the model distinguish tokens that have identical local context but
different positions in the sentence.

On top of the DeBERTa backbone, a standard sequence classification head
is added: a dropout layer followed by a linear projection from the
hidden dimension (768 for the base model) to 3 output logits, one for
each NLI label. During training, the \texttt{[CLS]} token's
representation --- which aggregates information from the entire input
sequence through the attention mechanism --- is passed through this
head to produce the final prediction.

%------------------------------------------------------------------
\subsection{Input Format}
%------------------------------------------------------------------

The input to Variant A is straightforward. The premise and hypothesis
are concatenated and passed to the tokenizer, which produces the
following token sequence:

\begin{center}
\texttt{[CLS] premise [SEP] hypothesis [SEP] [PAD] \ldots}
\end{center}

No explanation is included. This is the defining characteristic of
Variant A and the key contrast with Variants B, C, and D. The model
sees only the two sentences it must reason about, with no additional
context about why their relationship holds.

The maximum sequence length is set to 256 tokens. For Variant A this
is sufficient because premise and hypothesis together rarely exceed
50--70 tokens in e-SNLI. Choosing 256 rather than the 384 used in
Variants B, C, and D reduces memory usage and speeds up training
without any information loss, since no third text segment (the
explanation) needs to be accommodated.

%------------------------------------------------------------------
\subsection{Training Configuration}
%------------------------------------------------------------------

All hyperparameters follow the shared configuration defined in
\texttt{config.py} to ensure that performance differences between
variants can be attributed to architectural choices rather than
training differences.

\subsubsection{Label Mapping}

The label mapping is fixed across all variants:

\begin{verbatim}
label2id = {"entailment": 0,
            "neutral": 1,
            "contradiction": 2}
id2label = {0: "entailment",
            1: "neutral",
            2: "contradiction"}
\end{verbatim}

This exact ordering must be consistent across all variants. If Variant
A maps entailment to 0 but Variant B maps it to 2, the resulting
logits are incomparable and the evaluation numbers become meaningless.

\subsubsection{Sequence Length}

\begin{verbatim}
max_seq_length: int = 256
\end{verbatim}

Variant A uses 256 rather than the 384 used in Variants B, C, and D.
Since Variant A's input contains only the premise and hypothesis --- no
explanation --- sequences rarely approach even 100 tokens. Using 384
would waste memory and compute without any benefit.

\subsubsection{Training Hyperparameters}

\begin{verbatim}
learning_rate:    2e-5
num_train_epochs: 3
weight_decay:     0.01
warmup_steps:     3090
fp16:             True
\end{verbatim}

The learning rate of \texttt{2e-5} is the standard for fine-tuning
large pretrained transformers, established in the original BERT paper
and validated across hundreds of subsequent works. Too high (e.g.,
\texttt{1e-4}) destroys the pretrained representations through
catastrophic forgetting. Too low (e.g., \texttt{1e-6}) makes learning
impractically slow within three epochs.

Weight decay of \texttt{0.01} applies L2 regularization to the model
weights, penalising very large parameter values and acting as a soft
regulariser against overfitting. This is the standard value used with
the AdamW optimiser.

The warmup schedule linearly increases the learning rate from near-zero
to its full value over the first 3,090 steps --- approximately 6\% of
total training steps across 3 epochs with effective batch size 32. At
the very start of training, the randomly initialised classification
head produces highly noisy gradients. Without warmup, these noisy
updates at full learning rate would permanently damage the carefully
pretrained DeBERTa backbone weights before the head has had a chance to
stabilise.

FP16 (half-precision floating point) halves the memory footprint of all
stored tensors by switching from 32-bit to 16-bit floats. This is what
makes it feasible to train DeBERTa-v3-base at all on a T4 GPU with
16GB of VRAM.

\subsubsection{Batch Size and Gradient Accumulation}

\begin{verbatim}
per_device_train_batch_size: 8
gradient_accumulation_steps: 4
effective_batch_size:        32
\end{verbatim}

The effective batch size of 32 is achieved by accumulating gradients
over 4 mini-batches of 8 examples each before performing an optimiser
step. This is a standard technique for simulating a larger batch size
when GPU memory is insufficient to hold 32 examples simultaneously. All
variants use this same effective batch size to ensure training dynamics
are comparable.

%------------------------------------------------------------------
\subsection{Data Preprocessing}
%------------------------------------------------------------------

The e-SNLI training data is loaded from CSV files containing premise,
hypothesis, label, and up to three human-written explanations per
example. For Variant A, the explanation columns are ignored entirely
--- only the premise, hypothesis, and label are used.

Labels in the raw data are stored as integers (0, 1, 2). These are
mapped to their string equivalents (entailment, neutral, contradiction)
to match the label2id dictionary, and any rows with missing premise,
hypothesis, or label values are dropped. The training set is then
subsampled to exactly 100,000 examples using \texttt{random\_state=42}:

\begin{verbatim}
train_df = train_df.sample(
    n=100_000,
    random_state=42
)
\end{verbatim}

The fixed random seed is critical. Every variant must train on the
exact same 100K subset. If different variants train on different
subsets, any performance difference could reflect data variation rather
than architectural design, making the comparison invalid.

The final dataset sizes are shown in Table~\ref{tab:dataset_splits}.

\begin{table}[H]
\centering
\caption{Dataset splits used for Variant A training and evaluation.}
\label{tab:dataset_splits}
\begin{tabular}{lc}
\toprule
\textbf{Split} & \textbf{Examples} \\
\midrule
Training (subsampled) & 100,000 \\
Validation            & 9,842   \\
Test                  & 9,824   \\
\bottomrule
\end{tabular}
\end{table}

Label distribution in the training subset is approximately balanced:
33,523 entailment, 33,330 neutral, and 33,147 contradiction examples.

%------------------------------------------------------------------
\subsection{Training Procedure}
%------------------------------------------------------------------

Training follows a standard fine-tuning procedure using the
HuggingFace \texttt{Trainer} API. For each batch of 8 examples, the
tokenized premise-hypothesis pairs are passed through the DeBERTa
backbone. The \texttt{[CLS]} token representation from the final hidden
layer is projected through the classification head to produce 3 logits.
Cross-entropy loss is computed against the gold labels and
backpropagated through both the classification head and the backbone.
After every 4 mini-batches, the AdamW optimiser updates the model
weights.

At the end of each epoch, the model is evaluated on the validation set.
The checkpoint with the highest validation accuracy across all epochs is
retained as the best model, and this is the checkpoint used for final
test evaluation.

\begin{verbatim}
save_strategy          = "epoch"
save_total_limit       = 3
load_best_model_at_end = True
metric_for_best_model  = "accuracy"
\end{verbatim}

%------------------------------------------------------------------
\subsection{Training Results}
%------------------------------------------------------------------

Training was conducted on a Kaggle notebook with a T4 x2 GPU
configuration. Total wall time was approximately 3 hours and 20 minutes
across 3 epochs. Table~\ref{tab:training_results} summarises the loss
and accuracy at each epoch.

\begin{table}[H]
\centering
\caption{Training and validation metrics per epoch for Variant A.}
\label{tab:training_results}
\begin{tabular}{cccc}
\toprule
\textbf{Epoch} & \textbf{Train Loss} & \textbf{Val Loss} & \textbf{Val Acc} \\
\midrule
1 & 5.862 & 1.238 & 75.05\% \\
2 & 4.429 & 0.995 & 80.76\% \\
3 & 3.776 & 0.850 & \textbf{84.14\%} \\
\bottomrule
\end{tabular}
\end{table}

Both training and validation loss decreased consistently across all
three epochs. No overfitting was observed --- the validation loss
tracked the training loss throughout, suggesting the model was still
in an underfitting regime at epoch 3. Additional epochs would likely
yield further improvement, though compute constraints limited training
to three passes.

%------------------------------------------------------------------
\subsection{Evaluation Results}
%------------------------------------------------------------------

Following the evaluation protocol in \texttt{evaluate.py}, Variant A
was assessed on seven benchmarks spanning three categories. Results are
presented in Table~\ref{tab:eval_results}.

\begin{table}[H]
\centering
\caption{Variant A evaluation results across all benchmarks.}
\label{tab:eval_results}
\begin{tabular}{llc}
\toprule
\textbf{Benchmark}   & \textbf{Type}   & \textbf{Accuracy} \\
\midrule
e-SNLI Test          & In-domain       & 83.36\% \\
SNLI Test            & Cross-domain    & 83.36\% \\
MultiNLI Matched     & Cross-domain    & 58.45\% \\
MultiNLI Mismatched  & Cross-domain    & 61.08\% \\
ANLI R1              & Adversarial     & 31.00\% \\
ANLI R2              & Adversarial     & 31.50\% \\
ANLI R3              & Adversarial     & 32.00\% \\
\bottomrule
\end{tabular}
\end{table}

\subsubsection{In-domain Performance}

The 83.36\% test accuracy on e-SNLI is consistent with published
benchmarks for DeBERTa-v3-base on SNLI under similar data and compute
constraints. The gap between this and the state of the art
($\sim$91--92\%) is primarily attributable to the 100K subsampling ---
training on the full 550K examples would push accuracy significantly
higher. This subsampling is intentional and applied uniformly across
all variants to enable fair architectural comparison.

SNLI test accuracy matches e-SNLI test accuracy exactly at 83.36\%,
which is expected: both datasets were collected using the same
annotation methodology (Amazon Mechanical Turk workers writing
hypotheses for image captions), so they share the same vocabulary,
sentence structure, and label distribution.

\subsubsection{Cross-domain Generalization}

Performance drops substantially on MultiNLI: 58.45\% on the matched
set and 61.08\% on the mismatched set. This is a drop of roughly
22--25 percentage points from in-domain performance. MultiNLI spans
ten distinct genres of written and spoken English --- government
reports, telephone conversations, fiction, travel guides, and others
--- that are entirely absent from e-SNLI's image caption domain. The
model has learned to recognise entailment relationships in the specific
vocabulary and sentence patterns of image descriptions. When those
patterns disappear, so does much of its performance.

This cross-domain degradation is not a failure --- it is the
diagnostic that Variant A is designed to produce. It quantifies
precisely how much the model is relying on domain-specific shortcuts
rather than genuine linguistic reasoning.

\subsubsection{Adversarial Robustness}

The ANLI scores tell an even starker story. At 31.00\%, 31.50\%, and
32.00\% across rounds R1, R2, and R3 respectively, Variant A performs
at approximately random chance (33.3\% for a three-class problem). ANLI
was constructed through an adversarial human-in-the-loop process: human
annotators were specifically tasked with writing examples that would
fool a strong NLI model.

The fact that Variant A achieves near-random performance on ANLI is not
a failure of training --- it is a confirmation that the model has
learned exactly what we expected: surface patterns. When those patterns
are deliberately subverted, the model has nothing to fall back on. This
result is the strongest possible motivation for the
explanation-augmented approaches in Variants B, C, and D.

%------------------------------------------------------------------
\subsection{Role in the Project}
%------------------------------------------------------------------

Variant A establishes two things for the rest of the project. First,
it provides the numerical baseline: 83.36\% on e-SNLI test, 58--61\%
on MultiNLI, and $\sim$31--32\% on ANLI. Every other variant must be
evaluated against these numbers. An improvement on e-SNLI test
demonstrates that reasoning supervision helps in-domain. An improvement
on MultiNLI demonstrates better generalization. An improvement on ANLI
demonstrates genuine reasoning rather than pattern memorization.

Second, Variant A defines the lower bound for what is achievable
without reasoning. If any of the explanation-enhanced variants fail to
outperform Variant A on out-of-domain or adversarial benchmarks, that
would be a meaningful negative finding --- it would suggest that the
form of reasoning supervision used is not transferring to harder
evaluation conditions. The value of Variant A is not in its accuracy.
It is in what the accuracy reveals.

\end{document}